\section{Related Work}
\label{sec:related}

Carbon-aware scheduling has matured rapidly over the past four years.
We organise prior work into four threads: (i) batch and long-running
workloads, where the dominant approaches assume delay-tolerance and
suspend/resume; (ii) theoretical foundations of online carbon-aware
scheduling; (iii) the small but growing set of carbon-aware
schedulers explicitly targeting serverless or FaaS workloads, with
which \GreenFaaS{} most directly engages; and (iv) FaaS systems work
on cold-starts, warm pools, and characterization, on which our cost
model depends.

\subsection{Carbon-Aware Scheduling for Batch Workloads}

Wiesner et al.'s \emph{Let's Wait Awhile} \citep{wiesner2021lets}
introduced threshold-based deferral for batch workloads and released
the hourly carbon-intensity dataset (DE, GB, FR, US-CAISO, 2020) that
has become the de-facto reference; we use this dataset directly in
our evaluation (\S\ref{sec:evaluation:methodology}).
\emph{CarbonScaler} \citep{hanafy2023carbonscaler} scales the
parallelism of ML training jobs in response to carbon intensity,
exploiting the elasticity of data-parallel workloads.
Souza et al.'s \emph{GAIA} scheduler
\citep{sukprasert2023quantifying} addresses the three-way trade-off
between carbon, performance, and cost for batch jobs, and quantifies
the practical limits of carbon-aware spatiotemporal shifting; the same
group's \emph{Ecovisor} \citep{souza2023ecovisor} provides a virtual
energy system abstraction. Google's Carbon-Intelligent Compute
\citep{radovanovic2023google} deploys a Variable Capacity Curve
approach at production scale.
\emph{LACS} \citep{bostandoost2024lacs} is a learning-augmented
resource scaler addressing uncertain demand.
\emph{Adapting Datacenter Capacity} \citep{lin2023adapting} explored
provisioning to track grid carbon.

These works share two structural assumptions that \GreenFaaS{}
relaxes. \emph{First}, the schedulable units are \emph{jobs} with
cost separable per-job, allowing the schedulers to reason about each
job's deferral in isolation. FaaS invocations are not separable in
this sense because their cost is coupled through warm-pool state,
which we formalize in \S\ref{sec:problem}. \emph{Second}, the
schedulers operate on a timescale (minutes to hours) where cold-start
penalties are negligible. We show empirically in
\S\ref{sec:evaluation} that naively transplanting these techniques to
FaaS workloads --- as the Wait-Awhile baseline does ---
\emph{more than doubles} operational carbon relative to a
carbon-unaware FIFO baseline (163--236\% above FIFO across our
synthetic and real-carbon experiments). The ``batch schedulers don't
transfer'' claim is therefore not only intuitively obvious from the
assumption mismatch but also empirically severe.

\subsection{Theoretical Foundations}

A recent line of work has put carbon-aware online algorithms on
firmer footing. Lechowicz et al.'s \emph{Online Pause and Resume
Problem} \citep{lechowicz2024online} studies a one-dimensional model
in which a workload can be paused during high-carbon periods and
resumed during low-carbon periods, paying a fixed switching cost.
They derive double-threshold algorithms with provably optimal
competitive ratios. These theoretical results apply to a workload
model that is in several respects strictly simpler than ours: a
single workload (no multi-tenant mix), one machine (no spatial
routing), and a single switching cost (no idle-energy accumulation).
The FaaS setting we study introduces three structural features ---
per-invocation cold-start versus warm-pool trade-offs, spatial
routing across regions, and tiered SLA constraints --- that take us
outside the class of problems for which competitive-ratio bounds are
currently known. We are explicit in \S\ref{sec:problem:online} that
\GreenFaaS{} is an empirical scheduler with a \emph{local}
analytical guarantee (Lemma~\ref{lemma:tradeoff}'s break-even
characterization) rather than a competitively-optimal one. Extending
\citep{lechowicz2024online} to handle warm-pool idle energy and SLA
tiering is an open direction we identify in \S\ref{sec:discussion}.

\subsection{Carbon-Aware FaaS and Serverless}
\label{sec:related:faas}

Five recent systems target serverless workloads specifically; each
covers a different facet of the design space.

\paragraph{GreenCourier \citep{chadha2023greencourier}.} The closest
single-function spatial-routing precedent: a Kubernetes scheduler
plugin that routes serverless functions across geographically
distributed Knative clusters based on real-time grid carbon intensity.
It reports 9--18\% carbon reduction per invocation with negligible
latency overhead. GreenCourier and \GreenFaaS{} share the high-level
goal but differ in three significant respects.
\emph{First}, GreenCourier is \emph{spatial-only} --- it routes
functions to the cleanest available region, with no temporal deferral
and no warm-pool reasoning. Our headline contribution beyond
GreenCourier is therefore the joint spatial/temporal/warm-pool policy
and the underlying trade-off analysis.
\emph{Second}, GreenCourier has no analytical model of the warm-pool
versus cold-start carbon trade-off; our Lemma~\ref{lemma:tradeoff}
gives a closed-form characterization that becomes essential as soon
as the scheduler reasons about \emph{whether} to maintain a warm
pool.
\emph{Third}, GreenCourier's reported savings (9--18\%) sit at the
lower end of what our headline results suggest is achievable; we
attribute this to the spatial-only scope rather than to algorithmic
deficiencies in GreenCourier itself.

\paragraph{Caribou \citep{gsteiger2024caribou}.} Caribou is the
closest recent system in scope at the workflow level. It offloads
serverless \emph{workflows} (multi-function DAGs) across
geo-distributed regions for sustainability, automatically
(re-)deploying functions and redirecting traffic to lower-carbon
endpoints. Caribou requires no changes to application logic and
preserves performance and cost metrics. Its scope is workflow-level
spatial (no temporal deferral, no per-invocation warm-pool
reasoning), with decisions made at deployment time and refreshed by
periodic redeployment rather than per invocation.

We compare head-to-head in \S\ref{sec:evaluation:caribou} by porting
Caribou's deployment-solver logic to our per-invocation workload as
a single-node DAG case, with perfect carbon forecasting and constant
token-bucket grant (both favourable to Caribou). The finding:
Caribou(1h) and \GreenFaaS{} achieve carbon savings that differ by less than 0.1 percentage points
on our 24h FaaS slices; the two systems are operationally distinct
but empirically equivalent on per-invocation workloads with hour-scale
carbon variation. The systems are complementary along the granularity
axis: Caribou makes coarse-grained, deployment-time placement decisions
for whole workflows; \GreenFaaS{} makes fine-grained, per-invocation
decisions including warm-pool reasoning. A natural composition would
use Caribou for workflow placement and \GreenFaaS{} for per-invocation
scheduling within each chosen region; we have not implemented this
composition.

\paragraph{EcoLife \citep{jiang2024ecolife}.} EcoLife targets the
hardware-generation axis: it co-optimizes \emph{operational} and
\emph{embodied} carbon by exploiting multi-generation hardware (newer
servers are faster but have higher embodied carbon per unit time;
older servers are slower but already amortised). It uses a Particle
Swarm Optimisation-based scheduler that decides, per invocation, which
hardware generation to execute on. EcoLife explicitly reasons about
keep-alive (warm-pool) versus cold-start trade-offs --- which makes
it the closest precedent for our Lemma~\ref{lemma:tradeoff} --- but
it does so through a metaheuristic optimizer without a closed-form
characterization. The two contributions are largely orthogonal:
EcoLife answers ``on what hardware generation should this function
execute?''; \GreenFaaS{} answers ``in what region and at what time
should this function execute, and is a warm pool worth maintaining
there?''. A scheduler combining the two axes is a natural follow-up
(\S\ref{sec:discussion:embodied}).

\paragraph{CASA \citep{qi2024casa}.} CASA is the closest recent
single-cluster, autoscaling-focused competitor. It is an SLO- and
carbon-aware framework that schedules and autoscales containers in a
serverless cluster, explicitly trading off the SLO cost of cold-starts
against the carbon cost of pre-warming and keep-alive. CASA reports up
to a $2.6\times$ reduction in operational carbon while reducing SLO
violations by up to $1.4\times$ over a state-of-the-art baseline. The
follow-on work \citep{qi2024wastewater} extends the same multi-objective
framework to co-optimize wastewater alongside carbon and SLOs. CASA
and \GreenFaaS{} address overlapping concerns through different
mechanisms. CASA's contribution is a carbon-aware autoscaling and
container-placement policy in a single cluster, evaluated by simulation
against a state-of-the-art autoscaler. \GreenFaaS's contribution is a
closed-form trade-off characterization
(Lemma~\ref{lemma:tradeoff}) that yields a parameter-free per-invocation
decision rule, combined with a multi-region routing policy and a
formal do-no-harm property
(\S\ref{sec:evaluation:topology},
\S\ref{sec:evaluation:realtopology}). Where CASA establishes that
carbon-aware autoscaling is worthwhile, \GreenFaaS{} establishes the
analytical conditions under which warm-pool maintenance is itself
carbon-positive --- a question CASA's heuristics do not pose
explicitly. The two approaches could be composed at different
granularities: CASA at the container-pool autoscaling layer, \GreenFaaS{}
at the per-invocation routing-and-warm-pool-gate layer.

\paragraph{GreenWhisk \citep{greenwhisk}.} GreenWhisk is a
system-level redesign of Apache OpenWhisk to expose energy and carbon
state to load-balancing algorithms, supporting both grid-connected and
grid-isolated modes. Its contribution is infrastructural: an
energy-interface API and a carbon-aware load-balancer integrated into
a production-grade FaaS platform. The scheduling algorithm itself is
intentionally minimal in scope; the value is the open-platform
abstraction. \GreenFaaS's algorithmic and analytical contributions are
compatible with this abstraction: a production GreenWhisk deployment
could load Algorithm~\ref{alg:greenfaas} as its load-balancing policy
without changes to either side. GreenWhisk's two-mode design (grid
intermittency) is a deployment dimension our analysis does not address
and that an integration would have to handle separately.

\paragraph{CASPER \citep{souza2024casper}.} Adjacent rather than
directly competing: targets distributed web services (MediaWiki-class
applications) under latency SLOs, using a mixed-integer program to
jointly provision servers and load-balance requests across cloud
regions. It addresses spatial routing and provisioning with SLO
constraints but treats requests at the application level (request
rates, not per-invocation arrivals) and does not engage with
cold-start economics. We position \GreenFaaS{} as the per-invocation,
FaaS-native analogue of CASPER, with the cold-start trade-off model as
a key technical distinction.

\subsection{Serverless Cold-Starts and Characterisation}

Our cost model in \S\ref{sec:problem} and the SLA-tiered policy in
\S\ref{sec:algorithm:sla} depend on a substantial body of FaaS
systems work. \emph{Serverless in the Wild}
\citep{shahrad2020serverless} is the canonical characterization of
production Azure Functions workloads; it established that
invocations are dominated by a small number of high-frequency
``hot'' functions, that durations follow heavy-tailed distributions,
and that arrival processes are bursty and event-driven. Our workload
generator and the function popularity model (Zipf with $s = 1.2$)
follow this characterization directly. The follow-up
\emph{Azure Functions Invocation Trace 2021}
\citep{zhang2021faster,azuretrace} gives per-invocation records over
two weeks; we use its schema as the basis for our trace loader
(\S\ref{sec:evaluation:methodology}).

The cold-start problem itself has a rich literature.
\emph{Catalyzer} \citep{du2020catalyzer} and \emph{Firecracker}
\citep{agache2020firecracker} attack cold-starts at the system level
via micro-VM and snapshot techniques. \emph{FaaSCache}
\citep{fuerst2021faascache} frames warm-pool management as a caching
problem and derives keep-alive policies from caching theory. We
borrow the warm-pool / keep-alive abstraction directly, but add the
carbon dimension: idle energy at intensity $C_r$ in region $r$ is
itself a cost that the scheduler must trade off against cold-start
energy at potentially different $C$. This is the substance of our
Lemma~\ref{lemma:tradeoff}.

A parallel measurement-focused line of work makes it possible to
attribute per-invocation carbon footprints to FaaS workloads with
high accuracy. \emph{Accountable Carbon Footprints}
\citep{sharma2024accountable} develops a Shapley-value-based
disaggregation methodology that fairly attributes shared hardware and
software emissions to individual function invocations, achieving
$>$99\% accuracy on a wide range of workloads. This work is
methodologically orthogonal --- it measures, we schedule --- but a
production-grade carbon-aware scheduler should integrate with such a
measurement substrate to close the loop from decisions to observed
outcomes.

\paragraph{Summary of positioning.} To our knowledge, no prior work
has formalized the cold-start carbon trade-off in closed form,
demonstrated the temporal-axis idle-energy correction
(\S\ref{sec:tradeoff:idle}), or established a do-no-harm property
across topology, SLA-mix, forecast-accuracy, and carbon-variability
axes. These are the empirical and analytical novelties on which the
rest of the paper builds.
